\studySubsection{calc-06-differential-applications-mvt-equations-inequalities}{第 6 讲 一元函数微分学的应用（二）——中值定理、微分等式与微分不等式}

\begin{knowledgeBox}
教材页码：第 164 页起。

本讲用于整理罗尔定理、拉格朗日中值定理、柯西中值定理、泰勒公式、微分等式证明和微分不等式证明。新增题目时必须检查闭区间连续、开区间可导等使用条件。
\end{knowledgeBox}

\studySubsection{calc-06-auxiliary-function-construction}{辅助函数构造：从目标等式反推中值定理}

\knowledgeAnchor[MATH1-KN-CALC-0113]{辅助函数构造}
\noindent{\footnotesize\textbf{索引：} \knowledgeIndexRef[MATH1-KN-CALC-0113]{辅助函数构造}}\par

\begin{knowledgeBox}
\textbf{定位：这类题真正考什么？}

典型题面是：
\[
\text{证明存在 }\xi\in(a,b)\text{，使一个含 }f(\xi),f'(\xi)
\text{ 的等式成立。}
\]
表面上，未知点 \(\xi\) 没有办法直接求；本质上，题目要求把目标等式
改写为某个中值定理的结论。辅助函数不是凭灵感“猜”出来的，而是按
\[
\boxed{\text{目标式}\ \longrightarrow\ \text{移项为零}\
\longrightarrow\ \text{识别导数}\ \longrightarrow\ \text{制造端点关系}}
\]
反向设计出来的。

本节的前置是罗尔定理、拉格朗日中值定理和柯西中值定理；后续可迁移到
积分等式、微分不等式、泰勒余项和一阶线性微分方程的积分因子。学习顺序是：
先掌握统一原理，再识别五类常见导数结构，最后用例题练习方法选择。

\textbf{一、统一原理：把目标做成一个“可被迫为零”的导数}

把待证等式移项，记为
\[
E(\xi)=0.
\]
若能找到函数 \(F\) 和在 \((a,b)\) 内处处非零的因子 \(w\)，使
\[
F'(x)=w(x)E(x),
\]
同时满足
\[
F\in C[a,b]\cap D(a,b),\qquad F(a)=F(b),
\]
则由罗尔定理，存在 \(\xi\in(a,b)\) 使 \(F'(\xi)=0\)。又因
\(w(\xi)\ne0\)，所以 \(E(\xi)=0\)。

这揭示了构造时必须同时解决的三个问题：
\begin{enumerate}
  \item \textbf{导数问题：}怎样让 \(F'\) 含有目标式 \(E\)？
  \item \textbf{端点问题：}怎样让 \(F(a)=F(b)\)，或让 \(F\) 有足够多的零点？
  \item \textbf{合法性问题：}构造后是否真的连续、可导，分母和乘子是否安全？
\end{enumerate}

辅助函数通常不唯一。给 \(F\) 加常数或乘非零常数，不改变罗尔点；
在保持与 \(E=0\) 等价的前提下，也可以先给目标式乘一个处处非零的因子，
再识别导数。考场上应选导数最短、条件最容易核验的那个，而不是追求形式新奇。

\textbf{二、最小模型：拉格朗日中值定理本身就是一次构造}

若目标是
\[
f'(\xi)=\frac{f(b)-f(a)}{b-a},
\]
就从 \(f\) 中减去经过两端点的割线：
\[
F(x)=f(x)-f(a)-\frac{f(b)-f(a)}{b-a}(x-a).
\]
此时 \(F(a)=F(b)=0\)，而
\[
F'(x)=f'(x)-\frac{f(b)-f(a)}{b-a}.
\]
对 \(F\) 用罗尔定理，正好得到目标式。由此应形成一个稳定认识：
\[
\boxed{\text{端点不等时，减去插值函数，把“差”改造成“同”。}}
\]
两个端点减一次函数；三个节点减二次插值多项式；要得到高阶导数结论，
再对零点反复使用罗尔定理。

\textbf{三、五类常见构造及其识别信号}

\begin{enumerate}
  \item \textbf{移项后直接积分：}
  若目标是 \(f(\xi)=c\)，先看能否令
  \[
  F(x)=\int_a^x\bigl(f(t)-c\bigr)\,\mathrm dt.
  \]
  只要题设能保证 \(F(b)=F(a)=0\)，罗尔定理就给出
  \(F'(\xi)=f(\xi)-c=0\)。识别信号是“目标含函数值而不含导数”，
  且常数 \(c\) 由一个定积分给出。

  \item \textbf{乘指数因子：}
  看到
  \[
  f'(x)+p(x)f(x)
  \]
  时，若 \(p\) 在相关区间连续，取 \(P'(x)=p(x)\)，则
  \[
  \left(e^{P(x)}f(x)\right)'
  =e^{P(x)}\bigl(f'(x)+p(x)f(x)\bigr).
  \]
  因而辅助函数优先取 \(F=e^P f\)。若目标是 \(f'-pf\)，则改用
  \(F=e^{-P}f\)。指数因子始终为正，最后从 \(F'(\xi)=0\) 推回目标式
  不会丢失信息。

  \item \textbf{乘幂或除幂：}
  当 \(x^m\) 在整个区间有定义且不为零（特别是 \(x>0\)），
  且目标出现 \(xf'(x)-mf(x)\) 时，
  \[
  \left(\frac{f(x)}{x^m}\right)'
  =\frac{xf'(x)-mf(x)}{x^{m+1}},
  \]
  所以尝试 \(F=f/x^m\)。若出现 \(xf'(x)+mf(x)\)，则在定义合法时
  尝试 \(F=x^m f\)，因为
  \[
  (x^m f(x))'=x^{m-1}\bigl(xf'(x)+mf(x)\bigr).
  \]
  识别信号是“\(f'\) 前有 \(x\)，而 \(f\) 前是常数倍”。

  \item \textbf{乘积或商的导数：}
  看到 \(f'g+fg'\) 就想到 \((fg)'\)；看到
  \(f'g-fg'\) 就想到商法则：
  \[
  \left(\frac{f}{g}\right)'
  =\frac{f'g-fg'}{g^2},\qquad g\ne0.
  \]
  商式构造必须先检查 \(g\) 在整个闭区间上不为零，不能只检查目标点。

  \item \textbf{插值与多次罗尔降阶：}
  若题目给出三个或更多函数值，却要求二阶或更高阶导数在某点取指定值，
  先用多项式 \(p\) 拟合这些已知值，再令 \(F=f-p\)。若 \(F\) 有
  \(n+1\) 个互异零点，并具备相应阶数的可导性，则反复使用罗尔定理，
  可得到某点 \(F^{(n)}(\xi)=0\)，即
  \[
  f^{(n)}(\xi)=p^{(n)}(\xi).
  \]
\end{enumerate}

\textbf{补充：什么时候直接用柯西中值定理？}

若目标含两个导数及两个端点增量，可构造
\[
F(x)=\bigl(g(b)-g(a)\bigr)f(x)
     -\bigl(f(b)-f(a)\bigr)g(x).
\]
因为 \(F(b)-F(a)=0\)，罗尔定理给出
\[
\bigl(g(b)-g(a)\bigr)f'(\xi)
-\bigl(f(b)-f(a)\bigr)g'(\xi)=0.
\]
这正是柯西中值定理的交叉相乘形式，而且不必预先除以可能为零的
\(g'(\xi)\)。若题目最终需要比值，再单独核验分母不为零。

\textbf{四、方法选择：看到题面后怎样决定第一步？}

\begin{enumerate}
  \item 先把目标等式全部移到左边，得到 \(E(\xi)=0\)；不要一上来就写
  “令 \(F(x)=\cdots\)”。
  \item 从 \(E\) 中找导数指纹：和差看乘积、商；\(f'+pf\) 看指数因子；
  \(xf'-mf\) 看幂函数；只有函数值时看积分。
  \item 再看题设给了哪种端点信息：两个端点值适合罗尔或减割线；
  多个节点值适合插值加多次罗尔；两个函数的端点增量适合柯西中值定理。
  \item 若自然得到的 \(F\) 端点不相等，优先问“减去什么简单函数能让
  端点相等”，而不是立刻推翻原构造。
  \item 最后才誊写定理条件。构造正确但漏掉连续、可导、分母非零，
  证明仍不闭合。
\end{enumerate}

\textbf{例 1（基础模型：减去割线）}

\textbf{题面信号：}已知 \(f\in C[a,b]\cap D(a,b)\)，证明存在
\(\xi\in(a,b)\) 使
\[
f'(\xi)=\frac{f(b)-f(a)}{b-a}.
\]

\textbf{分析与选择：}目标是“导数减常数等于零”，而常数正是端点割线
斜率，因此减去割线函数。

\textbf{解：}令
\[
F(x)=f(x)-f(a)-\frac{f(b)-f(a)}{b-a}(x-a).
\]
由 \(a<b\) 及题设，\(F\in C[a,b]\cap D(a,b)\)，且
\[
F(a)=0,\qquad F(b)=f(b)-f(a)-\bigl(f(b)-f(a)\bigr)=0.
\]
由罗尔定理，存在 \(\xi\in(a,b)\) 使 \(F'(\xi)=0\)。又
\[
F'(x)=f'(x)-\frac{f(b)-f(a)}{b-a},
\]
故所求等式成立。

\textbf{易错点与核验：}不能漏掉 \(a<b\)，否则割线斜率无定义。取
\(f(x)=ux+v\) 检验时，两边都恒等于 \(u\)，说明符号与系数正确。

\textbf{例 2（指数因子：端点按指数比例给出）}

\textbf{题面信号：}设 \(f\in C[0,1]\cap D(0,1)\)，且
\[
f(1)=e^{-2}f(0).
\]
证明存在 \(\xi\in(0,1)\) 使
\[
f'(\xi)+2f(\xi)=0.
\]

\textbf{分析与选择：}目标左端是 \(f'+2f\)，它是
\(e^{2x}f(x)\) 的导数去掉正因子后的结果；端点比例恰好能让
\(e^{2x}f(x)\) 两端相等。

\textbf{解：}令 \(F(x)=e^{2x}f(x)\)。则
\[
F(0)=f(0),\qquad F(1)=e^2f(1)=f(0),
\]
且 \(F\in C[0,1]\cap D(0,1)\)。由罗尔定理，存在
\(\xi\in(0,1)\) 使
\[
0=F'(\xi)=e^{2\xi}\bigl(f'(\xi)+2f(\xi)\bigr).
\]
由于 \(e^{2\xi}>0\)，目标式成立。

\textbf{易错点与核验：}若误取 \(e^{-2x}f(x)\)，导数中会出现
\(f'-2f\)，符号相反。取 \(f(x)=e^{-2x}\) 时，题设和结论处处成立，
可快速核对指数符号。

\textbf{例 3（除以幂函数：把比例条件变成端点等值）}

\textbf{题面信号：}设 \(m\in\mathbb R\)，
\(f\in C[1,e]\cap D(1,e)\)，且
\[
f(e)=e^m f(1).
\]
证明存在 \(\xi\in(1,e)\) 使
\[
\xi f'(\xi)=m f(\xi).
\]

\textbf{分析与选择：}端点条件比较的是 \(f(x)\) 与 \(x^m\) 的增长，
目标又是 \(xf'-mf\)，故考察比值 \(f/x^m\)。

\textbf{解：}由于 \(x>0\)，实数幂 \(x^m\) 在 \([1,e]\) 上有定义且
不为零。令
\[
F(x)=\frac{f(x)}{x^m}.
\]
则
\[
F(1)=f(1),\qquad F(e)=\frac{f(e)}{e^m}=f(1).
\]
由罗尔定理，存在 \(\xi\in(1,e)\) 使 \(F'(\xi)=0\)。而
\[
F'(x)=\frac{xf'(x)-mf(x)}{x^{m+1}},
\]
分母在该区间不为零，故 \(\xi f'(\xi)-mf(\xi)=0\)。

\textbf{易错点与核验：}若区间经过 \(0\)，\(f/x^m\) 未必连续甚至未必
有定义，不能照搬。取 \(f(x)=Cx^m\) 时，\(F\) 为常数，目标式处处
成立。

\textbf{例 4（积分造原函数：从平均值反推零导数）}

\textbf{题面信号：}设 \(f\in C[a,b]\)，\(a<b\)。证明存在
\(\xi\in(a,b)\) 使
\[
f(\xi)=\frac{1}{b-a}\int_a^b f(t)\,\mathrm dt.
\]

\textbf{分析与选择：}目标只有 \(f(\xi)\)，没有导数；右侧是常数。
把 \(f-\text{平均值}\) 积分，就能让它成为某个辅助函数的导数。

\textbf{解：}记
\[
\bar f=\frac{1}{b-a}\int_a^b f(t)\,\mathrm dt,\qquad
F(x)=\int_a^x\bigl(f(t)-\bar f\bigr)\,\mathrm dt.
\]
由 \(f\) 连续及微积分基本定理，
\(F\in C[a,b]\cap D(a,b)\)，且
\[
F(a)=0,\qquad
F(b)=\int_a^b f(t)\,\mathrm dt-(b-a)\bar f=0.
\]
由罗尔定理，存在 \(\xi\in(a,b)\) 使
\[
0=F'(\xi)=f(\xi)-\bar f.
\]
结论成立。

\textbf{易错点与核验：}这里不能只写“令 \(F'=f-\bar f\)”而不说明
一个可用的原函数及端点关系。若 \(f\equiv c\)，则 \(\bar f=c\)，
结论对所有内点成立。

\textbf{例 5（插值加两次罗尔：把三个函数值降到二阶导数）}

\textbf{题面信号：}设 \(f\in C[-1,2]\)，在 \((-1,2)\) 内二阶可导，
且
\[
f(-1)=0,\qquad f(0)=1,\qquad f(2)=0.
\]
证明存在 \(\xi\in(-1,2)\) 使 \(f''(\xi)=-1\)。

\textbf{分析与选择：}题设给三个节点值，目标含二阶导数。先找通过这
三个点的二次多项式
\[
p(x)=\frac{-x^2+x+2}{2},
\]
再让 \(F=f-p\) 具有三个零点。

\textbf{解：}直接核验 \(p(-1)=0,p(0)=1,p(2)=0\)。令
\[
F(x)=f(x)-p(x).
\]
则 \(F(-1)=F(0)=F(2)=0\)。分别在 \([-1,0]\) 与 \([0,2]\)
使用罗尔定理，存在
\[
u\in(-1,0),\qquad v\in(0,2)
\]
使 \(F'(u)=F'(v)=0\)。再对 \(F'\) 在 \([u,v]\) 上使用罗尔定理，
存在 \(\xi\in(u,v)\subset(-1,2)\) 使
\[
0=F''(\xi)=f''(\xi)-p''(\xi)=f''(\xi)+1.
\]
故 \(f''(\xi)=-1\)。

\textbf{条件检查与核验：}三个零点必须互异；第二次使用罗尔定理需要
\(F'\) 在 \([u,v]\) 上连续、在 \((u,v)\) 内可导，这由 \(f\) 在
原开区间内二阶可导保证。取 \(f=p\) 时，\(f''\equiv-1\)，可核验
目标常数。

\textbf{五、失分防线：错误写法为什么错？}

\begin{enumerate}
  \item \textbf{只写“令某函数”，不交代动机。}
  正确表达应先写“目标移项为 \(E(\xi)=0\)，注意到
  \(w(x)E(x)\) 是某函数的导数”，再给出 \(F\)。这既说明构造来源，
  也便于检查符号。

  \item \textbf{有导数形式，却没有端点关系。}
  例如 \(F(x)=x\) 在 \([0,1]\) 上连续可导，但 \(F(0)\ne F(1)\)，
  所以不能由罗尔定理推出某点 \(F'=0\)。

  \item \textbf{端点相等，却漏掉可导性。}
  \(F(x)=|x|\) 在 \([-1,1]\) 上满足 \(F(-1)=F(1)\)，但在 \(0\)
  不可导；其所有可导点的导数只有 \(\pm1\)，不存在导数为零的点。

  \item \textbf{乘子可能为零，却强行约去。}
  在 \([-1,1]\) 上取 \(F(x)=x^2/2\)，则
  \[
  F'(x)=x=x\cdot1,\qquad F(-1)=F(1).
  \]
  罗尔点可以是 \(0\)，但不能从 \(x\cdot1=0\) 推出 \(1=0\)。
  因此由 \(F'=wE\) 推回 \(E=0\) 时，必须核验 \(w(\xi)\ne0\)。

  \item \textbf{商式或幂式跨过零点。}
  构造 \(f/g\) 必须保证 \(g\) 在整个相关区间不为零；构造
  \(f/x^m\) 时要检查 \(x^m\) 的定义与零点。不能因为最后的
  \(\xi\) “也许不在零点”就跳过闭区间上的连续性检查。

  \item \textbf{多次罗尔只数次数，不数零点。}
  要推出某点 \(F^{(n)}=0\)，通常需要至少 \(n+1\) 个互异零点，
  并逐层写出新零点所在的开区间；重复写同一个零点不增加次数。
\end{enumerate}

\textbf{六、考场可直接誊写的证明骨架}

\[
\begin{array}{l}
\text{将目标式移项为 }E(\xi)=0。\\
\text{注意到 }w(x)E(x)=\dfrac{\mathrm d}{\mathrm dx}F(x)，
\text{故令 }F(x)=\cdots。\\
\text{由题设可知 }F\in C[a,b]\cap D(a,b)，
\text{且 }F(a)=F(b)。\\
\text{由罗尔定理，存在 }\xi\in(a,b)\text{，使 }F'(\xi)=0。\\
\text{又 }F'(\xi)=w(\xi)E(\xi)，\ w(\xi)\ne0，
\text{故 }E(\xi)=0。
\end{array}
\]

若是高阶问题，把第四行改为：“\(F\) 有 \(n+1\) 个互异零点，
连续使用 \(n\) 次罗尔定理，得某点 \(F^{(n)}(\xi)=0\)。”

\textbf{七、掌握标准}

\begin{itemize}
  \item \textbf{基础过关：}看到给定的辅助函数，能完整核验连续、可导、
  端点相等，并能求对导数；能识别乘积、商、指数因子和幂函数四类结构。
  \item \textbf{130 分以上：}面对新的存在性等式，能先移项，再在两分钟
  内选出“积分、指数、幂、插值”中的首选路线；若端点不等，能主动减去
  一次函数；证明中不漏分母和乘子条件。
  \item \textbf{满分能力：}能处理含参数、退化情形和多个节点的高阶问题；
  能解释构造为何产生，能用最小反例指出某个错误证明究竟缺少哪一条条件，
  并能用特殊函数回代核验常数与符号。
\end{itemize}

\textbf{总结口令：}

\[
\boxed{\text{先把目标化零，再把它认成导数；两端做成相同，
条件逐项核实。}}
\]
最危险的三处是：导数符号写反、构造后端点并不相等、从
\(w(\xi)E(\xi)=0\) 推结论时忘记检查 \(w(\xi)\ne0\)。本节对应
“中值定理存在性证明”题型；下一步应把同一逆向设计法迁移到含积分、
高阶导数和微分不等式的陌生外观中。
\end{knowledgeBox}

\knowledgeAnchor[MATH1-KN-CALC-0060]{泰勒公式与常用展开}
\noindent{\footnotesize\textbf{索引：} \knowledgeIndexRef[MATH1-KN-CALC-0060]{泰勒公式与常用展开}}\par

\begin{knowledgeBox}
\textbf{一句话本质：}

泰勒公式是在展开点 \(a\) 附近，用一个多项式记录函数的局部信息。令
\[
T_n(x)=\sum_{k=0}^{n}\frac{f^{(k)}(a)}{k!}(x-a)^k,
\]
则 \(T_n\) 与 \(f\) 在 \(a\) 点的函数值、一阶导数、二阶导数，直到 \(n\) 阶导数全部相同：
\[
T_n^{(k)}(a)=f^{(k)}(a),\qquad k=0,1,\ldots,n.
\]
因此这些系数不是背出来的任意系数，而是被“逐阶匹配导数”唯一确定的。可以把 \(T_n\) 理解为 \(f\) 在 \(a\) 附近的 \(n\) 次局部模型。

\textbf{两种考研常用余项：}

若 \(f\) 在 \(a\) 点具有直到 \(n\) 阶的导数，则当 \(x\to a\) 时，带皮亚诺余项的泰勒公式为
\[
f(x)=\sum_{k=0}^{n}\frac{f^{(k)}(a)}{k!}(x-a)^k
+o\!\left((x-a)^n\right).
\]
它告诉我们误差比 \((x-a)^n\) 更高阶，最适合求极限、判阶和反求参数。

若 \(f\) 在连接 \(a\) 与 \(x\) 的闭区间上直到 \(n\) 阶连续可导，并在相应开区间内存在 \(n+1\) 阶导数，则存在介于 \(a\) 与 \(x\) 之间的 \(\xi\)，使
\[
f(x)=\sum_{k=0}^{n}\frac{f^{(k)}(a)}{k!}(x-a)^k
+\frac{f^{(n+1)}(\xi)}{(n+1)!}(x-a)^{n+1}.
\]
这是带拉格朗日余项的泰勒公式，最适合估计误差、判断余项符号和证明不等式。注意：\(n\) 次泰勒多项式的拉格朗日余项要用 \(n+1\) 阶导数。

\textbf{麦克劳林公式：}

展开点 \(a=0\) 时称为麦克劳林公式。当 \(x\to0\) 时，常用展开为
\[
\begin{aligned}
e^x&=1+x+\frac{x^2}{2!}+\frac{x^3}{3!}+\frac{x^4}{4!}+o(x^4),\\
\sin x&=x-\frac{x^3}{3!}+\frac{x^5}{5!}+o(x^5),\\
\cos x&=1-\frac{x^2}{2!}+\frac{x^4}{4!}+o(x^4),\\
\ln(1+x)&=x-\frac{x^2}{2}+\frac{x^3}{3}-\frac{x^4}{4}+o(x^4),\\
(1+x)^\alpha&=1+\alpha x+\frac{\alpha(\alpha-1)}{2!}x^2
+\frac{\alpha(\alpha-1)(\alpha-2)}{3!}x^3+o(x^3),\\
\frac{1}{1-x}&=1+x+x^2+x^3+o(x^3),\\
\arctan x&=x-\frac{x^3}{3}+\frac{x^5}{5}+o(x^5).
\end{aligned}
\]

\textbf{以 \(\ln(1+x)\) 为例：展开式是什么意思，怎样推出？}

式子
\[
\ln(1+x)=x-\frac{x^2}{2}+\frac{x^3}{3}+o(x^3),
\qquad x\to0,
\]
严格地说，是存在一个余项 \(r(x)\)，使
\[
\ln(1+x)=x-\frac{x^2}{2}+\frac{x^3}{3}+r(x),
\qquad
\lim_{x\to0}\frac{r(x)}{x^3}=0.
\]
因此 \(o(x^3)\) 表示“比 \(x^3\) 更高阶、更快趋于零”的量，并不表示它等于零。例如
\[
x^4=o(x^3),
\]
因为 \(x^4/x^3=x\to0\)；但 \(x^3\) 不是 \(o(x^3)\)，因为二者之比恒为 \(1\)。对于 \(\ln(1+x)\)，下一项实际上是 \(-x^4/4\)，它当然可以被收进 \(o(x^3)\) 中。

现在从泰勒公式推导。令
\[
f(x)=\ln(1+x).
\]
在展开点 \(0\) 处，
\[
\begin{aligned}
f(0)&=0,\\
f'(x)&=\frac{1}{1+x},& f'(0)&=1,\\
f''(x)&=-\frac{1}{(1+x)^2},& f''(0)&=-1,\\
f'''(x)&=\frac{2}{(1+x)^3},& f'''(0)&=2.
\end{aligned}
\]
代入三阶麦克劳林公式：
\[
\begin{aligned}
f(x)
&=f(0)+f'(0)x+\frac{f''(0)}{2!}x^2
  +\frac{f'''(0)}{3!}x^3+o(x^3)\\
&=0+x-\frac{1}{2}x^2+\frac{2}{6}x^3+o(x^3)\\
&=x-\frac{x^2}{2}+\frac{x^3}{3}+o(x^3).
\end{aligned}
\]
一般地，
\[
f^{(n)}(0)=(-1)^{n-1}(n-1)!,
\]
所以第 \(n\) 次项的系数为
\[
\frac{f^{(n)}(0)}{n!}=\frac{(-1)^{n-1}}{n},
\]
这就解释了系数为什么依次为 \(1,-\frac12,\frac13,-\frac14,\ldots\)。

不要把每个变式都单独背诵。像
\[
1-\cos x=\frac{x^2}{2}-\frac{x^4}{24}+o(x^4),
\qquad
x-\sin x=\frac{x^3}{6}+o(x^3)
\]
都应由基本展开现场相减得到。

\textbf{展开到哪一阶：}

先看最终要和谁比较，再决定阶数，而不是一律展开三项。
\begin{enumerate}
  \item 先令局部变量 \(h=x-a\to0\)，把所有函数写成关于 \(h\) 的形式。
  \item 若分母是 \(h^m\)，就把分子展开到足以找到相消后第一个非零项；要计算极限，通常至少要确定 \(h^m\) 的系数。
  \item 含参数时，先检查首个系数会不会因参数取值而变成零；若会，必须分情况并继续展开。
  \item 乘积、商式和复合式只保留会影响目标阶数的项，避免无效高阶计算。
\end{enumerate}
例如
\[
\begin{aligned}
\lim_{x\to0}\frac{\ln(1+x)-x+\frac12x^2}{x^3}
&=\lim_{x\to0}\frac{\left(x-\frac12x^2+\frac13x^3+o(x^3)\right)-x+\frac12x^2}{x^3}\\
&=\frac13.
\end{aligned}
\]
这里若只展开到二阶，恰好全部相消，仍然得不到答案；三阶项才是决定项。

\textbf{展开点与变量搬运：}

基本公式要求括号里的小量趋于 \(0\)。例如 \(x\to1\) 时令 \(h=x-1\)，则
\[
\ln x=\ln(1+h)=h-\frac{h^2}{2}+\frac{h^3}{3}+o(h^3).
\]
若 \(x\to\infty\)，常令 \(h=1/x\to0\) 后再展开。看到 \(e^{u(x)}\)、\(\ln(1+u(x))\) 或 \(\sin u(x)\) 时，也要先确认 \(u(x)\to0\)。

\textbf{常见误区：}
\begin{itemize}
  \item 把“多项式部分”直接写成与原函数恒等，却漏写余项；有限阶展开只是在展开点附近的近似。
  \item 在加减相消中只用一阶等价无穷小。例如“\(\sin x\sim x\)”不能直接算出 \(x-\sin x\) 的首项，必须展开到三阶。
  \item 展开阶数不足，或忽略参数会使首项系数为零；判断标准始终是“相消后第一个非零项”。
  \item 混淆泰勒公式与泰勒级数。有限阶公式带余项即可成立；把 \(n\to\infty\) 后的级数写成原函数，还必须证明余项趋于零。即使函数无穷次可导，也不自动等于其泰勒级数。
  \item 机械套麦克劳林公式而不检查展开点、定义域和小量是否趋于零。
\end{itemize}

\textbf{考场口令：}

“先定展开点，再看目标阶；逐阶相消，抓首个非零项；要估误差，就保留拉格朗日余项。”
\end{knowledgeBox}
